\documentclass[11pt]{article}
\usepackage[margin=1in]{geometry}
\usepackage[T1]{fontenc}
\usepackage[utf8]{inputenc}
\usepackage{lmodern}
\usepackage{amsmath,amssymb,amsthm}
\usepackage{graphicx}
\usepackage{booktabs}
\usepackage{microtype}
\usepackage{xcolor}
\usepackage[colorlinks=true,linkcolor=blue!55!black,citecolor=blue!55!black,urlcolor=blue!55!black]{hyperref}

\newtheorem{theorem}{Theorem}
\newtheorem{lemma}[theorem]{Lemma}
\newtheorem{corollary}[theorem]{Corollary}
\newtheorem{remark}[theorem]{Remark}
\newcommand{\D}{\mathbb D}
\newcommand{\T}{\mathbb T}
\newcommand{\N}{\mathbb N}
\newcommand{\Z}{\mathbb Z}
\newcommand{\Lp}{L^p}

\title{Iterated angular--radial limits settle the\\literal Bessel--Fourier convergence conjecture}
\author{Candidate full affirmative answer to Conjecture 1 in arXiv:2202.06119}
\date{13 August 2026}

\begin{document}
\maketitle

\begin{center}
\fcolorbox{black}{gray!8}{\parbox{0.91\linewidth}{\raggedright
\textbf{Status: candidate full affirmative result under the printed
quantifiers; semantic review important.}
For every $4/3<p<4$ and every $f\in L^p(\D)$, there are cofinal rectangular
cutoffs $(N_k,M_k)$ for which the Bessel--Fourier partial sums converge to
$f$ in $L^p(\D)$.  In fact, the angular cutoffs $M_k\to\infty$ may be fixed
in advance and the radial cutoffs may be required to exceed any prescribed
lower bounds.  The proof does \emph{not} produce one path that works for all
$f$ simultaneously; that stronger quantifier order is not stated in the
conjecture and remains open.}}
\end{center}

\section{Source conjecture and quantifier scope}

Let $\D=[0,1]\times\T$ in polar coordinates with measure
$r\,dr\,d\theta$.  Acosta Babb~\cite{source} writes the Bessel--Fourier
rectangular sums as
\begin{equation}\label{eq:rect}
 S_{N,M}f(r,\theta)
 =\sum_{m=-M}^{M}\sum_{n=1}^{N}
 a_{m,n}(f)J_{|m|}(j_n^{|m|}r)e^{2\pi i m\theta}.
\end{equation}
Conjecture~1, on arXiv PDF page~6, asks whether, for every
$4/3<p<4$ and $f\in L^p(\D)$,
\[
 \|S_{N_k,M_k}f-f\|_{L^p(\D)}\longrightarrow0
\]
``for some appropriate choice'' of $N_k,M_k\in\N$.

\begin{figure}[ht]
  \centering
  \includegraphics[width=\textwidth]{figures/open_problem_crop.png}
  \caption{Conjecture~1 in arXiv:2202.06119, PDF page~6.  The image contains
  the complete printed statement.}
\end{figure}

The order of quantifiers printed in the article (and repeated unchanged in
the author's 2024 thesis) is
\begin{equation}\label{eq:literal}
 \forall f\in L^p(\D)\quad\exists (N_k,M_k)_{k\ge1}.
\end{equation}
The discussion immediately after the conjecture emphasizes the difficulty of
uniform operator bounds, which may indicate that the stronger order
\begin{equation}\label{eq:stronger}
 \exists (N_k,M_k)_{k\ge1}\quad\forall f\in L^p(\D)
\end{equation}
was informally intended.  The theorem below proves \eqref{eq:literal}, with
two useful strengthenings, but deliberately does not claim
\eqref{eq:stronger}.

\section{Main theorem}

Let $A>0$ denote the absolute constant appearing in the Balodis--C\'ordoba
cutoff geometry quoted in the source.  It is included below to show that the
diagonal path can also obey $N_k\ge AM_k+1$; the proof itself does not need
this relation.

\begin{theorem}[Literal conjecture, with arbitrary lower bounds]\label{thm:main}
Let $4/3<p<4$, let $f\in L^p(\D)$, and let $(M_k)$ be any sequence of
positive integers with $M_k\to\infty$.  Given any positive integers $L_k$,
there are strictly increasing integers $N_k$ such that
\begin{equation}\label{eq:lower}
 N_k\ge\max\{L_k,AM_k+1\}
\end{equation}
and
\begin{equation}\label{eq:conv}
 \|S_{N_k,M_k}f-f\|_{L^p(\D)}\longrightarrow0.
\end{equation}
In particular, the affirmative assertion of Conjecture~1 holds exactly as
printed.
\end{theorem}

\begin{theorem}[Uniformity on compact sets]\label{thm:compact}
Under the hypotheses of Theorem~\ref{thm:main}, if
$K\subset L^p(\D)$ is norm compact, the cutoffs $N_k$ can be chosen so that
\[
 \sup_{f\in K}\|S_{N_k,M_k}f-f\|_{L^p(\D)}\longrightarrow0
\]
while still satisfying \eqref{eq:lower} and strict increase.
\end{theorem}

\begin{corollary}[One path for any countable family]\label{cor:countable}
For every countable set $\{f_1,f_2,\ldots\}\subset L^p(\D)$ and every
$M_k\to\infty$, one can choose a single cofinal sequence $(N_k,M_k)$ along
which $S_{N_k,M_k}f_j\to f_j$ for every $j$.
\end{corollary}

\section{Proof intuition}

A rectangle has two independent directions.  Freeze the angular width first.
Angular Fourier partial sums converge in $L^p$, so retaining the modes
$|m|\le M$ approximates $f$ when $M$ is large.  For a fixed $M$, only
finitely many radial functions $f_m$ remain.  The classical one-dimensional
Fourier--Bessel theorem says that, for each fixed order $|m|$, its radial
partial sums converge in exactly the range $4/3<p<4$.  One common radial
cutoff works for this finite list.  Choosing that cutoff after the angular
width is chosen gives an iterated limit; a diagonal selection turns it into
the requested sequence of rectangles.

This mechanism needs no operator bound uniform in both indices.  That is why
it proves the function-dependent quantifiers \eqref{eq:literal} while leaving
the whole-space uniform-path formulation \eqref{eq:stronger} untouched.

\section{The two one-parameter limits}

For $f\in L^p(\D)$, define its angular coefficients
\[
 f_m(r)=\int_{\T} f(r,\theta)e^{-2\pi i m\theta}\,d\theta
 \qquad(m\in\Z).
\]
H\"older's inequality on the probability space $\T$ gives
\begin{equation}\label{eq:coeffcontract}
 \|f_m\|_{L^p((0,1),r\,dr)}\le \|f\|_{L^p(\D)}.
\end{equation}
Let
\begin{equation}\label{eq:angular}
 P_Mf(r,\theta)=\sum_{|m|\le M}f_m(r)e^{2\pi i m\theta}.
\end{equation}

\begin{lemma}[Angular limit]\label{lem:angular}
For $1<p<\infty$,
\[
 \|P_Mf-f\|_{L^p(\D)}\longrightarrow0.
\]
Moreover, $\sup_M\|P_M\|_{L^p(\D)\to L^p(\D)}<\infty$.
\end{lemma}

\begin{proof}
The M.~Riesz theorem gives the uniform $L^p(\T)$ bound for ordinary Fourier
partial sums.  Apply it to $f(r,\cdot)$ and integrate the $p$th power against
$r\,dr$.  Convergence follows first for finite angular trigonometric
polynomials with radial $L^p$ coefficients and then for all $f$ by density
and the uniform bound.
\end{proof}

For a nonnegative integer $\nu$, write $B_N^{\nu}$ for the first $N$ terms
of the one-dimensional Fourier--Bessel expansion of order $\nu$ on
$L^p((0,1),r\,dr)$.  The classical mean-convergence theorem of
Wing and Benedek--Panzone~\cite{wing,benedek} says that, for each fixed
$\nu$,
\begin{equation}\label{eq:radialclassical}
 B_N^{\nu}g\longrightarrow g
 \quad\hbox{in }L^p((0,1),r\,dr),
 \qquad \frac43<p<4.
\end{equation}
This is precisely the fixed-order result invoked in the source article.

\begin{lemma}[Fixed angular window]\label{lem:fixedM}
If $4/3<p<4$ and $M$ is fixed, then
\[
 \|S_{N,M}f-P_Mf\|_{L^p(\D)}\longrightarrow0
 \qquad(N\to\infty).
\]
For this fixed $M$, the operators $S_{N,M}-P_M$ are uniformly bounded in
$N$ on $L^p(\D)$.
\end{lemma}

\begin{proof}
The definitions give the exact separation
\begin{equation}\label{eq:separate}
 S_{N,M}f-P_Mf
 =\sum_{|m|\le M}
   \bigl(B_N^{|m|}f_m-f_m\bigr)e^{2\pi i m\theta}.
\end{equation}
Since multiplication by $e^{2\pi i m\theta}$ embeds
$L^p((0,1),r\,dr)$ isometrically into $L^p(\D)$, the triangle inequality
and \eqref{eq:radialclassical} give
\begin{equation}\label{eq:finitesum}
 \|S_{N,M}f-P_Mf\|_{L^p(\D)}
 \le \sum_{|m|\le M}
 \|B_N^{|m|}f_m-f_m\|_{L^p((0,1),r\,dr)}\longrightarrow0.
\end{equation}
For every fixed order, pointwise strong convergence in
\eqref{eq:radialclassical} and Banach--Steinhaus imply
\[
 \sup_N\|B_N^{|m|}\|<\infty.
\]
Combine these finitely many bounds with
\eqref{eq:coeffcontract} and \eqref{eq:separate}.
\end{proof}

\section{Proofs of the main results}

\begin{proof}[Proof of Theorem~\ref{thm:main}]
For every $k$, Lemma~\ref{lem:fixedM} permits us to choose $N_k$ so large
that it is strictly larger than $N_{k-1}$, satisfies \eqref{eq:lower}, and
\begin{equation}\label{eq:diagerror}
 \|S_{N_k,M_k}f-P_{M_k}f\|_{L^p(\D)}<2^{-k}.
\end{equation}
Then
\[
 \|S_{N_k,M_k}f-f\|_p
 \le 2^{-k}+\|P_{M_k}f-f\|_p\longrightarrow0
\]
by Lemma~\ref{lem:angular}.
\end{proof}

\begin{proof}[Proof of Theorem~\ref{thm:compact}]
We use twice the standard fact that uniformly bounded operators converging
strongly converge uniformly on every norm-compact set.  Lemma~\ref{lem:angular}
therefore gives
\[
 \sup_{f\in K}\|P_{M_k}f-f\|_p\longrightarrow0.
\]
For each fixed $k$, Lemma~\ref{lem:fixedM} and its uniform boundedness
statement give
\[
 \sup_{f\in K}\|S_{N,M_k}f-P_{M_k}f\|_p\longrightarrow0
 \qquad(N\to\infty).
\]
Choose $N_k$ successively to make the latter supremum smaller than $2^{-k}$
and to meet the lower bounds.  The triangle inequality finishes the proof.
\end{proof}

\begin{proof}[Proof of Corollary~\ref{cor:countable}]
At stage $k$, apply the finite-set version of Theorem~\ref{thm:compact} to
$\{f_1,\ldots,f_k\}$ and choose $N_k$ so the radial error is below $2^{-k}$
for all these functions.  For each fixed $j$, the angular error tends to zero
as $k\to\infty$, and the radial bound applies whenever $k\ge j$.
\end{proof}

\section{Quantifier caveat and the negative clause}

The proof selects $N_k$ after seeing $f$ (or a compact/countable family).
It therefore does not show
\[
 \sup_k\|S_{N_k,M_k}\|_{L^p(\D)\to L^p(\D)}<\infty,
\]
which by Banach--Steinhaus is the natural requirement for one fixed path to
converge for every $f$.  The source's comment about the difficulty of uniform
bounds concerns exactly this stronger issue.

The sentence ``for $p$ outside this range, convergence fails in general'' is
also sensitive to the same semantics.  Classical mean-convergence results
show failure of the full fixed-order partial-sum sequence outside the open
range (including endpoint failure in the standard uniform-boundedness
sense).  They do not automatically exclude every $f$-dependent sparse
subsequence.  This packet proves the affirmative clause exactly as printed;
it does not make a new claim about that adaptive-subsequence interpretation
of the negative clause.

\section{Focused proof and upgrade audits}

Eight focused checks were made before promotion:
\begin{enumerate}
\item The arXiv source, 2023 published version, and 2024 author thesis were
checked; all retain the quantifier ``for every $f$, for some choice.''
\item The angular coefficient map was checked directly in
\eqref{eq:coeffcontract}, so no unjustified pointwise $L^p$ trace is used.
\item The angular limit was reduced to the scalar M.~Riesz theorem and product
density, avoiding any false Plancherel step for $p\ne2$.
\item The fixed-order radial theorem was checked against the classical
Fourier--Bessel literature and the source's own stated range.
\item The exact separation \eqref{eq:separate} was checked against the
normalization and order convention $J_m=J_{|m|}$ in the source.
\item The finite-mode triangle estimate \eqref{eq:finitesum} avoids any
constant uniform in the angular order; finiteness is enough.
\item The diagonal choice was strengthened to arbitrary prescribed $M_k$,
arbitrary radial lower bounds, compact families, and countable families.
\item A universal whole-space path was considered.  The argument stops
precisely at the missing uniform bound in both indices, so that stronger
problem is explicitly not claimed.
\end{enumerate}

\section{Novelty check, limitations, and review recommendation}

A bounded search on 13 August 2026 covered the exact title and conjecture
wording, ``iterated limit,'' ``diagonal,'' and ``fixed angular''
Bessel--Fourier searches, the 2023 publication, the author's 2024 thesis, the
citing triangular-domain paper, arXiv records through 2026, and the author's
2026 coefficient-decay paper.  The thesis still labels the problem open, and
no later solution or this diagonal argument was found.  Novelty is plausible,
not certified, and priority is not claimed.

The mathematical proof has no computational dependency.  Its limitation is
semantic rather than technical: it settles the literal
$\forall f\,\exists\text{ path}$ assertion, not the stronger
$\exists\text{ path}\,\forall f$ assertion that the surrounding discussion
may have intended.  Human review should first decide the intended quantifier
order.  Under the printed order, check only the fixed-window identity
\eqref{eq:separate} and the applicability of fixed-order mean convergence;
the rest is a finite-sum diagonal argument.

\begin{thebibliography}{9}
\bibitem{source}
R.~L. Acosta Babb,
\emph{Remarks on the $L^p$ convergence of Bessel--Fourier series on the disc},
Comptes Rendus Math\'ematique \textbf{361} (2023), 1075--1080;
arXiv:2202.06119, doi:10.5802/crmath.464.

\bibitem{wing}
G.~M. Wing,
\emph{The mean convergence of orthogonal series},
Amer. J. Math. \textbf{72} (1950), 792--808.

\bibitem{benedek}
A.~Benedek and R.~Panzone,
\emph{Mean convergence of series of Bessel functions},
Rev. Uni\'on Mat. Argentina \textbf{26} (1972/73), 42--61.
\end{thebibliography}

\end{document}
